This section aims to explain the main properties CRDT and OT in order to explain
their relevance in eventual consistent systems and compare them later on.

\subsection{CRDTs}

\subsubsection{Intuition}

In many cases, Data needs to be stored in copies on multiple computers also
known as replicas. Examples of these cases are \cite{CRDT_Overview}:

\begin{itemize}
    \item Mobile apps
    \item Distributed databases
    \item Collaboration software
    \item Large-scale data storages
\end{itemize}

All of those systems need to deal with the fact that data might be updated
concurrently on multiple devices.

\subsection{CRDT properties}

All CRDTs have these properties in order to ensure strong eventual consistency:

\begin{itemize}
    \item Associativity: The order of merging does not matter.
    \item Commutativity: Whoever does a merge with who has no influence.
    \item Idempotency: Merging the same change multiple times has no additional effect.
\end{itemize}

\subsection{State-based vs. Operation based}

\subsubsection{State-based CRDT}
State based CRDTs rely on sending their full state periodically to other
replicas in order to synchronise the system. This approach tends to become costly
as CRDT states grow. Delta-based CRDTs tackle this problem
by producing small incremental states (deltas) to be used for synchronization. Surprisingly,
current synchronisation algorithms induce redundant, wasteful deltas, which causes it to
perform worse than expected \cite{Stateful_delta}.

\subsubsection{Operation Based CRDT}

Changes in Operation based CRDT happen locally first. After that, the peer shares
the operation to other peers, which also operate locally. In order to ensure conflict-free
synchronisation between peers, these three conditions need to be satisfied \cite{operation_based_overview}:

\begin{itemize}
    \item Messages are received in the same order as sent.
    \item Every message is sent exactly once
    \item Operations, that can be processed simultaneously, have to be commutative
\end{itemize}

The drawback of this approach are the constraints it puts on the communication channel:
Messages have to be delivered only once (fire-and-forget), in causal order, to each peer\cite{sending_order}.

\subsection{Important Data Types}

\subsubsection{G-Counter}

The G-Counter is a state-based CRDT designed for a cluster of $n$ nodes. Each node is assigned a unique ID (ranging from $0$ to $n-1$)
and may only increment its own slot $P[i]$ in a shared array $P$. To retrieve the total count, all slots are summed. When two nodes synchronize,
their states are merged by taking the maximum value for each slot: $Z.P[i] = \max(X.P[i], Y.P[i])$. This ensures that no increments are lost,
regardless of the order in which updates arrive \cite{shapiro2011crdt}.


\subsubsection{PN-Counter}

As previously mentioned, the G-Counter revolves around incrementation only, but what if it has to be decremented?
This is what PN-Counter does. In contrast to G-Counter, it has two Counters, one tracking the increments and the other
tracking the decrements. Before merging, it subtracts the positive count with the negative one, then proceeds to merge
like a G-Counter, by comparing both values \cite{PN_Counter}.


\subsection{Operational Transformation}

\subsubsection{Intuition}

Operational Transformation is a conflict resolution algorithm in distributed systems
with more than one node accepting data changes. From a development perspective, collaborative
editing is a difficult task to approach with a lot of different solutions with their own trade-offs.
In order to imitate immediate response, each client stores its own local copy, otherwise clients would have to
wait for their actions to sync with other clients, which would cause a noticeable delay.
So the main challenge is to ensure consistency over peers which means, that they have to converge to a single version.

\subsubsection{Transformation Function}

Two concurrent Operations are converted in a way, that ensures, that the operations are commutative.
For example, if User A inserts an "X" at position 3 and User B deletes position 3 concurrently, it would
result in a conflict, but this is exactly where the transformation function helps. It transforms both
operations in a way that ensures, that there is no conflict.

\paragraph{How it works}

Before diving into transformation logic, it is crucial to understand how edits are represented.
Peers can not just send inserted text to the server, instead, they need to send an atomic operation
that describes following points:

\begin{itemize}
    \item what changed
    \item where it changed
    \item when it changed
\end{itemize}

In practice, the package the peer sends would look like this: First, it sends the DocId, which describes
the document the changes are assigned to. Then, it would send the base version (the version the user saw).
Afterward, it sends the userId representing the user who made the change and most important, the actual operation,
which is either an insert or a delete.

You can see an example of the json file in listing \ref{lst:msg_json}

\begin{lstlisting}[language=json, label=lst:msg_json]
{
    "docId": "doc_2345",
    "userId": "user_12",
    "baseVersion": 14,
    "op": {
        "type": "insert",
        "pos": 67,
        "text": "I"
    }
}
\end{lstlisting}
In order to help the server keep track, the client needs to store the current
version number of the document and the text state of the current document.

The server accepts the request and compares the base-version the client sent with
its own base version. In a simple scenario, the versions match and the operation can
operate without any transformation. However, if the server has a different state, the server
has to transform the operation. In this case, the server has saved all the operations made, so it
is able to fetch the old version of the document. Finally, it transforms the operation of the client
and transforms it against each historical operation in the order they were applied. As a result, the operation
can be applied without any conflict.

\paragraph{Convergence}

The core guarantee of OT is that all replicas converge to the same state once all operations have been
delivered and applied. Two properties are required to achieve this \cite{ellis1989concurrency}:

\begin{itemize}
    \item \textbf{TP1 (Transformation Property 1):} If two operations $O_1$ and $O_2$ are concurrent,
    then applying $O_1$ followed by $T(O_2, O_1)$ must yield the same document state as applying $O_2$
    followed by $T(O_1, O_2)$.
    \item \textbf{TP2 (Transformation Property 2):} Required in decentralized settings where operations
    can be transformed against other already-transformed operations. This property ensures the
    transformation function remains consistent across different transformation paths.
\end{itemize}

In a centralized architecture, only TP1 is required, which considerably simplifies the implementation.
TP2 is notoriously difficult to satisfy and remains an open research challenge for peer-to-peer OT systems.

\subsubsection{Limitations of OT}

While OT is well-suited for collaborative text editing, it comes with notable limitations:

\begin{itemize}
    \item \textbf{Complexity:} Designing a correct transformation function that satisfies TP1 and TP2 is
    non-trivial. Several published OT algorithms were later found to contain correctness bugs
    \cite{ellis1989concurrency}.
    \item \textbf{Centralized dependency:} Most practical OT implementations rely on a central server
    to serialize and order operations. This introduces a single point of failure and limits scalability
    in fully decentralized scenarios.
    \item \textbf{State overhead:} The server must retain the full operation history to be able to
    transform late-arriving operations from clients that are far behind the current document version.
\end{itemize}

Despite these drawbacks, OT remains widely deployed in production systems such as Google Docs,
where the centralized model is an acceptable architectural trade-off.

